\chapter{Preparação do Ambiente}

\section{Requisitos}

O projeto exige uma distribuição LaTeX atualizada com LuaLaTeX e Biber. Também é necessário um editor de texto com suporte a arquivos UTF-8 e, para conferência, um leitor de PDF.

Ambientes locais podem usar MiKTeX ou TeX Live. Em serviços online, selecione LuaLaTeX como compilador e envie a estrutura completa do projeto, preservando nomes e caminhos relativos.

\section{Obtenção da versão correta}

Obtenha o pacote pelo repositório controlado do projeto ou por canal institucional indicado. Não copie apenas \texttt{tcc.tex}: a classe, os módulos de \texttt{src/}, a identidade visual e os arquivos bibliográficos são necessários.

Depois de extrair ou clonar o projeto, confirme a presença de:

\begin{lstlisting}
undftex.cls
tcc.tex
src/
identidade/
obrigatorio/
opcional/
capitulos/
\end{lstlisting}

\section{Teste inicial}

Compile o modelo em branco antes de inserir conteúdo. O teste separa problemas de instalação de erros introduzidos durante a edição.

\begin{lstlisting}
lualatex -interaction=nonstopmode -halt-on-error tcc.tex
biber tcc
lualatex -interaction=nonstopmode -halt-on-error tcc.tex
lualatex -interaction=nonstopmode -halt-on-error tcc.tex
\end{lstlisting}

O resultado esperado é \texttt{tcc.pdf}. Mensagens sobre campos ainda não preenchidos são normais no primeiro teste; interrupção do processo, arquivo ausente ou comando desconhecido exigem correção.

\section{Editor e codificação}

Mantenha os arquivos em UTF-8. Não substitua caracteres acentuados por sequências incompatíveis nem altere a codificação de apenas um capítulo. Configure o editor para mostrar final de linha, pares de chaves e erros de sintaxe.
